\chapter{The C++ API}
\label{ch:cpp}

The entire engine layer is usable directly from C++17 --- the Python API is a
thin binding over it, so every concept of Chapters~\ref{ch:sa}--%
\ref{ch:optimal} maps one-to-one.

\section{Umbrella include and headers}

\begin{cppcode}
#include "qanneal/core.hpp"
\end{cppcode}

\begin{center}
\small
\begin{tabular}{@{}ll@{}}
\toprule
\textbf{Class / function} & \textbf{Header}\\
\midrule
\code{DenseIsing} & \code{dense\_ising.hpp}\\
\code{SparseIsing}, \code{SparseEdge} & \code{sparse\_ising.hpp}\\
\code{QUBO} & \code{qubo.hpp}\\
\code{AnnealSchedule} & \code{schedule.hpp}\\
\code{SQASchedule} & \code{sqa\_schedule.hpp}\\
\code{Annealer} & \code{annealer.hpp}\\
\code{ReplicaAnnealer} & \code{replica\_annealer.hpp}\\
\code{ParallelTemperingAnnealer} & \code{parallel\_tempering.hpp}\\
\code{SQAAnnealer} & \code{sqa\_annealer.hpp}\\
\code{SQAParallelTemperingAnnealer} & \code{sqa\_parallel\_tempering.hpp}\\
\code{CTPIMCAnnealer} & \code{ctpimc\_annealer.hpp}\\
\code{State}, \code{SQAState} & \code{state.hpp}, \code{sqa\_state.hpp}\\
Observers (all) & \code{metrics\_observer.hpp}\\
\code{MPIContext}, \code{run\_replica\_anneal()} &
\code{mpi/mpi\_context.hpp}, \code{mpi/mpi\_replica.hpp}\\
\bottomrule
\end{tabular}
\end{center}

\section{A complete C++ example}

Dense spin glass, simulated annealing:

\begin{cppcode}
#include "qanneal/core.hpp"
#include <iostream>
#include <random>
#include <vector>

int main() {
    const std::size_t n = 64;
    std::mt19937_64 rng(7);
    std::normal_distribution<double> g(0.0, 1.0);

    std::vector<double> h(n, 0.0);
    std::vector<std::vector<double>> J(n, std::vector<double>(n, 0.0));
    for (std::size_t i = 0; i < n; ++i)
        for (std::size_t j = i + 1; j < n; ++j)
            J[i][j] = J[j][i] = g(rng);

    qanneal::DenseIsing ising(h, J, /*c=*/0.0);
    auto schedule = qanneal::AnnealSchedule::linear(0.1, 5.0, 80);

    qanneal::Annealer annealer(ising, schedule);
    annealer.set_seed(42);
    auto result = annealer.run(/*sweeps_per_beta=*/60);

    std::cout << "best energy: " << result.best_energy << "\n";
    for (auto s : result.best_state.spins) std::cout << int(s) << ' ';
    std::cout << std::endl;
    return 0;
}
\end{cppcode}

And the quantum analogue in a few lines:

\begin{cppcode}
auto sqa_sched = qanneal::SQASchedule::from_vectors(betas, gammas);
qanneal::SQAAnnealer sqa(ising, sqa_sched,
                         /*trotter_slices=*/32, /*replicas=*/4);
auto res = sqa.run(/*sweeps_per_beta=*/40,
                   /*worldline_sweeps=*/4,
                   /*cluster_sweeps=*/1);
// adaptive schedule:
auto opt = sqa.run_optimal(/*beta=*/1.0,
                           /*j_perp_start=*/0.1, /*j_perp_end=*/25.0,
                           /*eps_tilde=*/0.0,      // budget calibration
                           /*alpha=*/15.0/14.0,
                           /*num_steps=*/150, /*sweeps_per_step=*/20,
                           /*worldline_sweeps=*/3);
\end{cppcode}

Reference build targets and CMake flags are in Chapter~\ref{ch:install};
\code{examples/dense\_ising\_sa.cpp},
\code{examples/sqa\_quantum\_parallel\_tempering.cpp} and
\code{examples/mpi/sa\_mpi.cpp} in the repository are compiled by the
default build and serve as living documentation.

\section{Implementation notes for contributors}

\begin{itemize}
\item \textbf{Local-field caches.} Every engine keeps
$f_i=h_i+\sum_jJ_{ij}s_j$ per (replica, slice); proposals are $O(1)$
(Eq.~\eqref{eq:deltaE}), accepted flips update the cache along the flipped
spin's neighbourhood. Backends implement
\code{energy()}, \code{delta\_energy()} and
\code{compute\_local\_fields()} over raw \code{int8\_t*} slices.
\item \textbf{Threading model.} OpenMP parallelises over replicas; each
replica owns disjoint regions of the field cache, its own
\code{std::mt19937\_64} stream seeded from the master RNG, its own shuffled
visit order and cluster buffers. No locks are needed in the hot path.
\item \textbf{Bond-sum tracking.} The adaptive schedule's $B$
(Eq.~\eqref{eq:chiB}) is updated incrementally by every move
(Section~\ref{sec:chib-est}) --- keep this invariant when adding new move
types: a move must report its exact $\Delta B$.
\item \textbf{Observers.} A sweep-level observer forces the sequential SQA
path (deterministic interleaving for reproducible traces); metrics-only
observers keep the parallel path.
\item \textbf{Errors.} Parameter validation throws
\code{std::invalid\_argument} (surfaced as \code{ValueError} in Python);
sentinel resolutions that alter semantics print a \code{cerr} warning rather
than failing silently.
\end{itemize}
